% !TeX root=../main.tex

\chapter{مقدمه}
\label{chap:introduction}

\section{نرم‌افزار در تعامل با جهان}

بخش مهمی از نرم‌افزارهای امروز فقط ورودی ثابتی را دریافت نمی‌کنند تا پس از محاسبه خاتمه یابند. یک کنترل‌کنندهٔ آسانسور باید در تمام مدت کار ساختمان به درخواست طبقه، بازشدن در، مانع میان در و وضعیت موتور واکنش نشان دهد. یک خدمت‌دهندهٔ شبکه باید هم‌زمان پیام چند کاربر، پایان مهلت‌ها و قطع ارتباط را مدیریت کند. چراغ راهنمایی، کنترل‌کنندهٔ صنعتی و سامانهٔ پایش بیمار نیز پیوسته داده‌ای از محیط می‌گیرند و رفتار بعدی خود را بر اساس حالت جاری انتخاب می‌کنند. این دسته از نرم‌افزارها نمونه‌هایی از \glspl{ReactiveSystem} هستند؛ یعنی سامانه‌هایی که معنای رفتارشان در یک تعامل ادامه‌دار با محیط شکل می‌گیرد \cite{SirjaniEtAl2004Rebeca,Sirjani2004Thesis}.

در یک برنامهٔ ترتیبی کوچک، یافتن خروجی نادرست برای یک ورودی معمولاً مسیر نسبتاً روشنی برای رفع خطا فراهم می‌کند. در سامانهٔ واکنشی، نتیجه به تاریخچهٔ رویدادها نیز وابسته است. ممکن است دو پیام جداگانه درست پردازش شوند، اما رسیدن آن‌ها با ترتیبی خاص به بن‌بست بینجامد؛ یا خطا تنها پس از تکرار طولانی یک چرخه آشکار شود. اگر چند مؤلفه مستقل و ناهمگام باشند، تعداد ترتیب‌های ممکن اجرای همان چند رویداد نیز به‌سرعت افزایش می‌یابد. بنابراین مسئله فقط درست‌بودن یک تابع نیست، بلکه باید رفتار همهٔ اجراهای مرتبط بررسی شود.

زمان در برخی از این سامانه‌ها جزئی از تعریف درستی است. در پخش صوت، دیررسیدن یک بسته شاید تنها به افت کیفیت منجر شود؛ اما در ترمز خودرو، تشخیص برخورد یا رهگیری یک هدف، پاسخ درست پس از پایان مهلت دیگر پاسخ قابل قبول نیست. چنین سامانه‌ای \gls{RealTimeSystem} نامیده می‌شود. صفت «بلادرنگ» الزاماً به معنای سریع‌بودن مطلق نیست؛ معنای دقیق‌تر آن این است که رویداد یا پاسخ باید در بازهٔ زمانی تعیین‌شده رخ دهد. برای نمونه، کنترل دمای یک اتاق می‌تواند مهلتی چندثانیه‌ای داشته باشد، درحالی‌که کنترل یک فرایند فیزیکی سریع به مهلتی بسیار کوتاه‌تر نیاز دارد. آنچه اهمیت دارد پیوند رفتار منطقی و محدودیت زمانی است \cite{AlurDill1990,AlurDill1994}.

\section{هزینهٔ یک فرض زمانی نادرست}

رخداد سامانهٔ دفاع موشکی پاتریوت در ظهران نمونه‌ای روشن از اهمیت همین پیوند است. در ۲۵ فوریهٔ ۱۹۹۱، سامانه نتوانست یک موشک اسکاد ورودی را رهگیری کند و برخورد موشک با محل استقرار نیروها به کشته‌شدن ۲۸ نفر انجامید. بررسی دفتر حسابرسی عمومی ایالات متحده نشان داد محاسبهٔ زمان در رایانهٔ کنترل تسلیحات بر اثر تبدیل عددی دقت از دست می‌داد. خطای کوچک هر محاسبه با طول مدت کار سامانه انباشته می‌شد؛ سامانه در زمان حادثه بیش از ۱۰۰ ساعت پیوسته کار کرده بود و در نتیجه، محدودهٔ جست‌وجوی هدف را در محل نادرستی محاسبه می‌کرد \cite{GAO1992Patriot}.

این رخداد تنها یک «خطای گردکردن» جدا از طراحی سامانه نبود. پاتریوت در ابتدا برای استقرارهای کوتاه طراحی شده بود و استفادهٔ پیوستهٔ طولانی در فرض‌های اولیهٔ آن قرار نداشت. رفتار نرم‌افزار برای چند ساعت می‌توانست پذیرفتنی باشد، اما همان محاسبه در افق زمانی طولانی ویژگی لازم برای رهگیری را نقض می‌کرد. یک مشخصهٔ صوری که کران خطای زمان یا قرارداشتن هدف در پنجرهٔ رهگیری را برای کل مدت مجاز کار بیان کند، این فرض پنهان را به یک تعهد قابل بررسی تبدیل می‌کند. این گزاره به معنای آن نیست که هر استفاده از روش صوری حتماً حادثه را پیش‌بینی می‌کرد؛ اگر محدودهٔ کار یا ویژگی مورد انتظار اشتباه مشخص شود، اثبات نیز همان فرض اشتباه را معتبر می‌گیرد. ارزش روش صوری در واداشتن طراح به بیان دقیق چنین فرض‌هایی و بررسی همهٔ حالت‌های مدل‌شده است.

شکست نخستین پرواز آریان ۵ در سال ۱۹۹۶ جنبهٔ دیگری از مسئله را نشان می‌دهد. موشک حدود ۴۰ ثانیه پس از آغاز پرواز از مسیر خارج شد و از بین رفت. گزارش رسمی آژانس فضایی اروپا علت را خطاهای مشخصه و طراحی در نرم‌افزار سامانهٔ مرجع اینرسی دانست؛ دو واحد افزونه نیز به‌دلیل اجرای نرم‌افزار یکسان تقریباً هم‌زمان از کار افتادند. آزمون‌های انجام‌شده مسیر واقعی آریان ۵ را به‌اندازهٔ کافی بازنمایی نکرده بودند و در نتیجه پیش‌شرط نامعتبر تبدیل عددی کشف نشد \cite{ESA1996Ariane}. این مثال یادآوری می‌کند که افزونگی سخت‌افزاری در برابر خطای مشترک طراحی محافظت کافی ایجاد نمی‌کند و پوشش آزمون نیز به سناریوهای انتخاب‌شده محدود است.

\section{از آزمون تا درستی‌یابی صوری}

آزمون نرم‌افزار همچنان بخش ضروری فرایند مهندسی است: اجرای واقعی می‌تواند خطای پیاده‌سازی، ناسازگاری محیط و مسئله‌ای را آشکار کند که در مدل انتزاعی حضور ندارد. محدودیت بنیادی آزمون این است که هر مورد آزمون تنها یک یا چند اجرای مشخص را مشاهده می‌کند. در سامانه‌های هم‌روند و زمان‌دار، فضای ترتیب پیام‌ها، انتخاب‌ها و مقدارهای زمانی ممکن است چنان بزرگ باشد که پیمایش دستی یا اجرای همهٔ آزمون‌ها عملی نباشد.

در \gls{FormalVerification}، مدل سامانه و ویژگی مورد انتظار با زبان دقیق ریاضی بیان می‌شوند. سپس یک روش استدلالی یا الگوریتمی بررسی می‌کند آیا ویژگی برای رفتارهای مدل برقرار است یا خیر. \gls{ModelChecking} یکی از این روش‌هاست که فضای حالت مدل را به‌طور خودکار جست‌وجو می‌کند و در صورت نقض ویژگی معمولاً یک مثال نقض ارائه می‌دهد \cite{ClarkeGrumbergPeled1999,BaierKatoen2008}. تضمین حاصل نه دربارهٔ هر جزئیات جهان واقعی، بلکه دربارهٔ مدل، مشخصه و فرض‌های تعریف‌شده است؛ ازاین‌رو، آزمون و درستی‌یابی صوری مکمل یکدیگرند.

کاربرد روش‌های صوری به پژوهش دانشگاهی محدود نمانده است. مهندسان \lr{Amazon Web Services} از سال ۲۰۱۱ مشخصه‌نویسی صوری و وارسی مدل با \lr{TLA+} را برای مسائل دشوار طراحی سامانه‌های توزیع‌شدهٔ بحرانی به کار گرفته‌اند \cite{NewcombeEtAl2015}. در Meta، تحلیل‌گر ایستای \lr{Infer} با استدلال نمادین و منطق ریاضی، مسیرهای متعدد برنامه را بدون اجرای آن بررسی می‌کند و در زمان گزارش رسمی، هر ماه صدها خطای احتمالی پیش از انتشار رفع می‌شدند؛ البته این ابزار تضمین کامل درستی تابعی ارائه نمی‌کند \cite{MetaInfer2015}. پروژهٔ \lr{Everest} در مایکروسافت نیز مؤلفه‌های رمزنگاری، تجزیه‌گرها و پروتکل‌های ارتباط امن را همراه با اثبات‌های ماشین‌بررسی‌شده برای استفادهٔ عملی توسعه داده است \cite{MicrosoftEverest2019}. این نمونه‌ها طیفی از روش‌های صوری را نشان می‌دهند: از تحلیل خودکار برای یافتن رده‌ای از خطاها تا اثبات ویژگی‌های مشخص یک مؤلفه.

\section{مسئلهٔ انفجار فضای حالت}

پوشش همهٔ رفتارهای یک مدل هزینه دارد. اگر هر مؤلفه چند حالت داشته باشد، ترکیب هم‌زمان چند مؤلفه می‌تواند تعداد بسیار بیشتری حالت سراسری بسازد. صف پیام، مقدار متغیرها، ترتیب اجرای کنشگرها و وضعیت ساعت‌ها نیز این فضا را گسترش می‌دهند. این پدیده \gls{StateSpaceExplosion} نام دارد و یکی از موانع اصلی کاربرد \gls{ModelChecking} است \cite{ClarkeGrumbergPeled1999,BaierKatoen2008}.

بااین‌حال، همهٔ جزئیات مدل برای هر پرسش تحلیلی اهمیت یکسان ندارند. فرض کنید هدف، بررسی توالی فرمان‌های قابل مشاهدهٔ یک کنترل‌کننده باشد. گام‌های داخلی محاسبه یا پیام‌هایی که در سطح انتزاع مورد نظر دیده نمی‌شوند، ممکن است مسیر اجرا را طولانی کنند بی‌آنکه مشاهدهٔ بیرونی تازه‌ای بسازند. اگر بتوان حالت‌هایی را که از دید رفتار مورد توجه تمایزپذیر نیستند شناسایی و ادغام کرد، تحلیل روی مدلی کوچک‌تر انجام می‌شود. شرط اساسی آن است که کاهش، همان جنبهٔ رفتاری را که قرار است بررسی شود حفظ کند.

در سامانهٔ زمان‌دار، صرفاً حذف کنش داخلی کافی نیست. دو اجرا ممکن است از نظر توالی پیام‌های قابل مشاهده یکسان باشند، اما مدت متفاوتی داشته باشند. شبیه‌سازی دوسویهٔ ضعیف زمان‌دار اجازه می‌دهد کنش‌های داخلی نادیده گرفته شوند، درحالی‌که رفتار قابل مشاهده و مقدار زمان سپری‌شده حفظ می‌شود \cite{ReynierSangnier2009,LiebEtAl2025}. رده‌های هم‌ارزی حاصل می‌توانند مبنای ساخت یک مدل خارج‌قسمتی کوچک‌تر باشند.

\section{بستر ربکا و افرا}

\gls{RebecaLanguage} زبانی کنشگرمحور برای مدل‌سازی سامانه‌های هم‌روند و توزیع‌شده است. هر شیء واکنشی حالت و صف پیام خود را دارد و از طریق ارسال ناهمگام پیام با دیگر اشیا ارتباط برقرار می‌کند. \gls{TimedRebecaLanguage} سازوکارهایی برای تأخیر، مهلت و زمان محاسبه به این مدل می‌افزاید \cite{Sirjani2004Thesis,SirjaniEtAl2004Rebeca,ReynissonEtAl2014TimedRebeca}. این ساختار برای بیان شهودی یک سامانهٔ واکنشی مناسب است، اما ترکیب صف‌ها، کنشگرها و زمان همان مسئلهٔ رشد فضای حالت را ایجاد می‌کند.

\gls{AfraTool} محیط مدل‌سازی و تحلیل خانوادهٔ ربکاست و با گردآوردن کامپایلر، وارسی‌گر مدل و ابزارهای نمایش، فضای حالت حاصل را برای تحلیل در اختیار کاربر می‌گذارد \cite{AfraWebsite2026,AfraGitHub2026}. ورودی پروژهٔ حاضر کد منبع ربکا نیست، بلکه سامانهٔ انتقال زمان‌داری است که این زنجیره با قالب \lr{.statespace} تولید می‌کند. بنابراین مسئله در مرز میان تولید مدل معنایی و تحلیل آن قرار دارد: چگونه می‌توان با انتخاب تعامل‌های قابل مشاهده، مدل را بر اساس شبیه‌سازی دوسویهٔ ضعیف زمان‌دار کاهش داد و هم‌زمان شاهدی قابل وارسی برای حفظ رفتار ارائه کرد؟

\section{هدف‌ها و پرسش‌های پروژه}

هدف اصلی این پروژه تبدیل تعریف و شبه‌کد کاهش ضعیف زمان‌دار به ابزاری اجرایی برای فضای حالت افراست. این پایان‌نامه الگوریتم هم‌ارزی تازه‌ای پیشنهاد نمی‌کند؛ کار آن تطبیق تعریف موجود با قرارداد واقعی خروجی افرا، پیاده‌سازی فرایند کاهش و فراهم‌کردن شواهد قابل تکرار برای درستی خروجی است.

پروژه سه پرسش را دنبال می‌کند:

\begin{enumerate}
\item آیا می‌توان خروجی زمان‌دار افرا را بدون ویرایش دستی خواند و به یک نمایش مستقل از رابط کاربری افرا تبدیل کرد؟
\item آیا می‌توان پس از پنهان‌سازی تعامل‌های انتخاب‌نشده، یک مدل خارج‌قسمتی کوچک‌تر ساخت و رابطهٔ آن با ورودی را به‌طور مستقل وارسی کرد؟
\item تصمیم معنایی دربارهٔ جمع‌شدن تأخیرهای جداشده با کنش داخلی چه اثری بر هم‌ارزی نمونه‌ها دارد و روش روی ورودی‌های موجود چه مقدار کاهش ایجاد می‌کند؟
\end{enumerate}

برای پاسخ به این پرسش‌ها، هدف‌های اجرایی زیر تعیین شده‌اند:

\begin{itemize}
\item استخراج و مستندسازی بخش مورد نیاز از قرارداد فایل \lr{TTS} افرا؛
\item پیاده‌سازی پنهان‌سازی، پالایش تأخیر، ساخت گذار ضعیف و \gls{PartitionRefinement}؛
\item ساخت و فشرده‌سازی \gls{QuotientSystem} همراه با نگاشت حالت‌های ورودی؛
\item توسعهٔ وارسی‌گری مستقل برای آزمودن رابطهٔ مدل تحویلی و ورودی؛
\item ثبت مصنوعات، سنجه‌ها و آزمون‌های لازم برای تکرار ارزیابی.
\end{itemize}

\section{دامنه و دستاوردهای پروژه}

دامنهٔ پیاده‌سازی به سامانه‌های انتقال متناهی با زمان گسسته و تأخیر صحیح نامنفی محدود است. ابزار خروجی \lr{TTS} افرا را می‌خواند و از \lr{FTTS}، زمان چگال، کامپایل مستقیم مدل ربکا و ادغام به‌صورت افزونهٔ رابط کاربری افرا پشتیبانی نمی‌کند. معنای پیش‌فرض، تأخیر چندواحدی را به گام‌های واحد می‌شکند تا تأخیرهای جداشده با کنش داخلی جمع‌پذیر باشند؛ خوانش سخت‌گیرانهٔ شبه‌کد نیز برای مقایسه حفظ شده است. فصل \ref{chap:method} این تصمیم را دقیق‌تر توضیح می‌دهد.

دستاورد پروژه یک ابزار مستقل جاوا با رابط خط فرمان، خوانندهٔ مقاوم در برابر ویژگی خاص فایل واقعی افرا، خط لولهٔ کاهش قابل ردیابی و وارسی‌گر جداگانه است. نتیجهٔ اصلی به‌صورت \lr{JSON} بدون اتلاف، نگاشت افراز، نمایش سازگار با افرا، نمودار و سنجه‌های اجرا ذخیره می‌شود \cite{AfraWeakTimedReduction2026}. ارزش این کار در پیوندزدن تعریف نظری با قرارداد داده و نیازهای مهندسی زنجیرهٔ ربکا و افراست؛ ادعای مینیمال یا کانونی‌بودن مدل خروجی و نیز ادعای مقیاس‌پذیری صنعتی مطرح نمی‌شود.

\section{روش کلی ارزیابی}

پیاده‌سازی در سه سطح بررسی شده است. آزمون‌های واحد، اجزای خواندن فایل، معنای گذار ضعیف، پالایش افراز، ساخت خارج‌قسمت و وارسی رابطه را می‌سنجند. آزمون‌های انتهابه‌انتها برنامهٔ بسته‌بندی‌شده را از بیرون اجرا و مصنوعات، کدهای خروج و قطعی‌بودن نتیجه را بررسی می‌کنند. در نهایت، چهار ورودی شامل یک نمونهٔ کوچک و سه بیان از مطالعهٔ موردی خانهٔ هوشمند ارزیابی شده‌اند. دو نمونهٔ خانهٔ هوشمند، تأخیر ده‌واحدی را به‌ترتیب به $7+\tau+3$ و $3+\tau+7$ می‌شکنند تا حساسیت روش به معنای زمان آشکار شود. سنجه‌های فصل \ref{chap:results} مستقیماً از خروجی ثبت‌شدهٔ همین اجراها گزارش می‌شوند.

\section{ساختار پایان‌نامه}

فصل \ref{chap:background} پیش‌نیازهای درستی‌یابی صوری، وارسی مدل، سامانه‌های واکنشی و زمان‌دار، ربکا و افرا، سامانه‌های انتقال و تعریف شبیه‌سازی دوسویهٔ ضعیف زمان‌دار را ارائه می‌کند. فصل \ref{chap:method} قرارداد ورودی، تصمیم‌های معنایی، مراحل الگوریتم، معماری پیاده‌سازی و روش اعتبارسنجی را شرح می‌دهد. فصل \ref{chap:results} نتیجهٔ کاهش، حساسیت معنایی، پوشش آزمون‌ها، زمان اجرا و تحلیل نتایج را گزارش می‌کند. در پایان، فصل \ref{chap:conclusion} به پرسش‌های پروژه پاسخ می‌دهد، دستاوردها و محدودیت‌ها را جمع‌بندی می‌کند و مسیرهای ادامهٔ کار را پیشنهاد می‌دهد.
